\documentclass[12pt,fontset=none]{ctexart}
\usepackage[a4paper,margin=1in]{geometry}
\usepackage{enumitem}
\usepackage{hyperref}
\usepackage{longtable}
\usepackage{array}
\usepackage{fancyhdr}
\usepackage{xcolor}
\setlength{\headheight}{15pt}
\setCJKmainfont{STSong}
\setCJKsansfont{STSong}
\setCJKmonofont{STSong}
\pagestyle{fancy}
\fancyhf{}
\rhead{Cell Biology Review}
\lhead{整理自 notebook.md 与课程 PDF}
\cfoot{\thepage}

\setlist[enumerate]{leftmargin=*, itemsep=0.8em}

\title{Cell Biology 问答整理}
\author{}
\date{2026-04-26}

\begin{document}
\maketitle
\tableofcontents
\newpage

\section*{说明}
本整理以当前目录中的 \texttt{notebook.md} 为主，并结合 \texttt{CB-Ch4-041626.pdf} 与 \texttt{CB-Ch67-042326.pdf} 的课程内容补全。少数题目存在术语拼写或表述歧义时，按细胞生物学常用写法统一为标准术语，例如 \texttt{Dynesin} 统一作 \texttt{dynein}，\texttt{t-tubulin} 按上下文理解为 \texttt{T-tubule}。

\addcontentsline{toc}{section}{细胞骨架}
\section*{细胞骨架}
\begin{enumerate}
\item \textbf{stress fibers是哪一种cytoskeleton?}：属于 \textbf{actin cytoskeleton}，本质上是非肌肉细胞中的 \textbf{contractile actin bundles}，常与 myosin II 协同产生张力。

\item \textbf{Contractile ring是哪一种cytoskeleton？}：属于 \textbf{actin cytoskeleton}；由 actin filament 和 myosin II 组成，在胞质分裂时收缩。

\item \textbf{neutrophil的移动和shape的改变依靠的是哪一种cytoskeleton}：主要依靠 \textbf{actin cytoskeleton}，尤其是皮层 actin network 的快速重组。

\item \textbf{请指出epithelial cells 的apical和basal}：\textbf{apical} 面朝向腔面或外界，常可见 microvilli；\textbf{basal} 面朝向 basal lamina / extracellular matrix，通过 hemidesmosomes 等与基底膜连接。

\item \textbf{将FtsZ、MreB、ParM与Crescentin进行一一对应}：
FtsZ 对应 \textbf{tubulin-like}，参与细菌胞质分裂并形成 Z-ring；
MreB 对应 \textbf{actin-like}，参与维持细胞形状；
ParM 对应 \textbf{actin-like}，参与质粒或染色体分离；
Crescentin 对应 \textbf{intermediate filament-like}，其 coiled-coil 特征类似中间丝。

\item \textbf{g-actin的三种亚型以及其主要作用都是什么？}：脊椎动物常讲的三类是 \textbf{$\alpha$-actin、$\beta$-actin、$\gamma$-actin}。$\alpha$-actin 主要在肌细胞中，负责收缩；$\beta$-actin 多见于非肌细胞，偏向细胞迁移、前沿延伸和胞质分裂；$\gamma$-actin 也常见于非肌细胞，更偏向皮层 actin 网络的稳定、细胞形态维持和张力调节。

\item \textbf{分辨nocodazole和Taxol的作用}：\textbf{nocodazole} 促进微管去聚合，抑制纺锤体形成；\textbf{Taxol} 稳定微管、抑制其动态不稳定性，阻止正常纺锤体重排。二者都会干扰有丝分裂，因此都偏向杀伤分裂细胞。

\item \textbf{cofilin和capZ 的作用}：\textbf{cofilin} 结合 actin filament 侧面并增加扭曲，削弱亚基间作用，促进切断和去聚合；\textbf{CapZ} 是 plus-end capping protein，封住 actin 丝正端，抑制其继续加长或缩短，从而稳定该端。

\item \textbf{actin filament的bundle和network形态分别是由什么蛋白介导的？}：\textbf{network} 主要由 \textbf{Arp2/3 complex} 介导分支成核；\textbf{bundle} 主要由 \textbf{formin} 促进形成长而直的 actin filament，再由 bundling/cross-linking proteins 组织成束。

\item \textbf{Parallel bundle 和 contractile bundle的区别是取决于什么以及是有什么形成的？}：关键取决于 \textbf{actin filament 的相对极性、丝间距以及交联蛋白种类}。\textbf{parallel bundle} 中各丝方向一致，常由 fimbrin/fascin 等紧密交联形成，如 microvilli；\textbf{contractile bundle} 中 actin 常为反向平行排列，丝间距较大，允许 \textbf{myosin II} 插入，常由 $\alpha$-actinin 等形成，如 stress fibers 与 contractile ring。

\item \textbf{Myosin II的结构是什么？}：myosin II 是由 \textbf{2 条 heavy chains 和 2 对 light chains} 构成的二聚体。每条 heavy chain 有一个头部 motor domain、颈部 light-chain-binding 区和长尾部 coiled-coil；多个 myosin II 可进一步装配成 \textbf{bipolar thick filament}。

\item \textbf{肌节的结构是什么？}：肌节（sarcomere）是两条 \textbf{Z disc} 之间的重复单位。薄丝 actin 的 plus end 锚定在 Z disc；粗丝 myosin II 位于中央；中央有 \textbf{M line} 连接粗丝；\textbf{titin} 提供弹性并帮助定位粗丝；\textbf{nebulin} 与 \textbf{tropomodulin} 参与薄丝长度控制。

\item \textbf{平滑肌的相应Ca2+信号进而产生收缩的机理是什么？}：平滑肌中 Ca$^{2+}$ 升高后不靠 troponin，而是先与 \textbf{calmodulin} 结合，激活 \textbf{MLCK}（myosin light-chain kinase），后者磷酸化 myosin regulatory light chain，使 myosin II 活化并与 actin 作用产生收缩；去磷酸化则使其松弛。

\item \textbf{$\alpha$-tubulin 和 beta-tubulin本质的差别是什么？}：二者共同形成 heterodimer，但 \textbf{$\alpha$-tubulin 上结合的 GTP 基本不水解、也不交换}；\textbf{$\beta$-tubulin 上的 GTP 可被水解并可交换}，因此微管动态不稳定性的关键化学变化发生在 $\beta$-tubulin。

\item \textbf{如果microtubules丢掉了GTP cap会发生什么？}：会发生 \textbf{catastrophe}。失去 GTP cap 后，GDP-tubulin protofilament 倾向弯曲外翻，微管迅速去聚合。

\item \textbf{$\gamma$-tubulin的功能是什么}：主要负责 \textbf{微管成核}。它富集在 MTOC/centrosome，组成 $\gamma$-tubulin 复合体，为新微管提供成核模板，常与微管 minus end 相关。

\item \textbf{Tubulin-Sequestering and Microtubule-Severing Proteins功能分别是什么？}：\textbf{tubulin-sequestering proteins} 结合游离 tubulin dimer，降低可用于聚合的 tubulin 浓度；\textbf{microtubule-severing proteins}（如 katanin）切断已有微管，重塑微管网络，同时也会产生更易去聚合的新末端。

\item \textbf{kinesin和dyneins的方向分别是什么？}：大多数 \textbf{kinesin} 朝 \textbf{plus end} 运动；大多数 \textbf{dynein} 朝 \textbf{minus end} 运动。

\item \textbf{Dynein attach to membrane enclosed organelles依靠的蛋白是什么？}：主要依靠 \textbf{dynactin} 介导 dynein 与膜性细胞器连接。

\item \textbf{Lissencephaly的名字是什么，为什么会导致这个？}：\textbf{lissencephaly} 中文常译为“\textbf{无脑回畸形}”或“平滑脑”。某类 lissencephaly 与 \textbf{Lis1} 缺陷有关；Lis1 是 dynein 结合蛋白，参与神经元发育时的核迁移。核迁移失败会导致皮层神经元定位异常，因此大脑皮层褶皱发育异常。

\item \textbf{Dynein、linker protein导致鞭毛波动的基本机制？}：轴丝中 dynein 让相邻微管双联体尝试彼此滑动；但 \textbf{linker proteins} 将双联体彼此约束，滑动被转化为 \textbf{弯曲}，从而形成纤毛/鞭毛的波动。

\item \textbf{Plectin、SUN-KASH的作用是什么？}：\textbf{Plectin} 属于 plakin 家族，把 intermediate filament 与微管、actin 或膜连接起来，起交联和力学整合作用；\textbf{SUN-KASH complex} 横跨核膜，把核纤层/核内骨架与细胞质骨架连接起来，实现核定位、机械传力和核迁移。

\item \textbf{Rac-GTP和Rho-GTP的功能分别是什么？}：\textbf{Rac-GTP} 主要促进 lamellipodia 和 branched actin network 形成，利于细胞前沿延伸；\textbf{Rho-GTP} 主要促进 stress fiber 和 focal adhesion 形成，增强收缩与张力。
\end{enumerate}

\addcontentsline{toc}{section}{细胞膜与膜转运}
\section*{细胞膜与膜转运}
\begin{enumerate}
\item \textbf{Cell membrane function是什么？（一共有四个）}：四个核心功能是 \textbf{选择性屏障}、\textbf{小分子转运}、\textbf{信号接收与传递}、\textbf{提供流动性以支持生长/运动/形变}。

\item \textbf{细胞膜主要由哪三大成分组成？}：主要由 \textbf{lipids、proteins、carbohydrates} 组成；若只按膜主体讲，通常强调脂质双层和膜蛋白，糖链多以糖脂和糖蛋白形式位于膜外侧。

\item \textbf{Phospholipid 里面最主要的是哪个成分？}：最丰富的磷脂通常是 \textbf{phosphatidylcholine}。

\item \textbf{胆固醇的主要作用是什么，fill 了因为什么而形成的gap?}：胆固醇主要 \textbf{稳定并调节膜流动性}，还能增加膜刚性；它填补的是 \textbf{不饱和脂肪酸尾部 cis-double bond 形成弯折后产生的空隙}。

\item \textbf{膜脂在脂双层中的运动模式以及频率？}：主要有四类。\textbf{lateral diffusion} 很快，可在约 1 秒内侧向移动约 2 $\mu$m；\textbf{rotation} 也很快，可达每秒数万转量级；\textbf{flexion} 最快，常在纳秒级；\textbf{flip-flop} 在无酶帮助时极少发生，可低到每月不到一次。

\item \textbf{membrane proteins有哪四种？}：四类常见功能型膜蛋白是 \textbf{transporters、anchors、receptors、enzymes}。

\item \textbf{参与信号蛋白膜定位的两类主要脂质锚是什么？}：主要是 \textbf{fatty acid anchors} 和 \textbf{prenyl groups}。

\item \textbf{水通道蛋白的结构是怎样的？}：aquaporin 常形成 \textbf{tetramer}，但每个单体本身就是一个独立水孔；每个单体通常含 \textbf{6 条跨膜 $\alpha$-helix}，以及 \textbf{2 条只跨半层膜的短螺旋}，中央有 NPA motif 构成选择性滤过区。

\item \textbf{transmembrane domains 的结构一般是什么？}：最常见是 \textbf{$\alpha$-helix}；较少见的是 \textbf{$\beta$-barrel}，常见于 porin 等通道蛋白。

\item \textbf{hydropathy plots怎么看？}：看序列疏水性随位置的变化。持续较高的 \textbf{正峰} 通常提示一段疏水跨膜 $\alpha$-helix；负值或低值区域多为亲水、位于膜外或胞质侧。峰值既要高，也要有足够长度，才更像真正跨膜区。

\item \textbf{glycosylated的位置是哪里}：膜蛋白的糖基化位点位于 \textbf{non-cytosolic side}，也就是 ER/Golgi 腔面，最终对应细胞表面的 \textbf{胞外侧}。

\item \textbf{nanodisc的意义和形成过程是怎样的？}：nanodisc 是一种 \textbf{人工小型脂双层平台}，用于在近天然膜环境中稳定和研究膜蛋白。通常先用 detergent 把膜蛋白和磷脂一起溶出，再加入 membrane scaffold protein，去除 detergent 后，脂双层自组装成圆盘，膜蛋白嵌入其中。

\item \textbf{什么蛋白必须需要detergent，而什么蛋白不需要？}：\textbf{integral membrane proteins} 需要 detergent 才能从膜中溶出；\textbf{peripheral membrane proteins} 通常不需要 detergent，用高盐、pH 改变或温和方法即可释放。

\item \textbf{Membrane-bending Proteins三种模式}：三种经典模式是 \textbf{hydrophobic insertion}、\textbf{scaffolding}、\textbf{protein crowding}。

\item \textbf{Permeases是什么？}：permeases 指 \textbf{transporters/carriers}，通过构象变化把溶质运过膜，不像通道那样形成持续开放的孔。

\item \textbf{transporter具有三种状态是什么？}：\textbf{outward-open、occluded、inward-open}。

\item \textbf{active transport的三种途径分别是什么？}：\textbf{coupled transport}、\textbf{ATP-driven pumps}、\textbf{light-driven pumps}。

\item \textbf{inverted repeats的意义是什么？}：许多 transporter 由 \textbf{倒置重复结构} 构成，形成伪对称性，便于通过两侧构象互换实现 \textbf{alternating access}；也提示它们常由基因复制和进化重排而来。

\item \textbf{Anti-diabetic与Anti-depressant}：这里对应 Na$^+$-coupled transporter 的药理学例子。\textbf{anti-diabetic} 常指 \textbf{SGLT2 inhibitors}，抑制肾脏 Na$^+$-glucose 共转运；\textbf{anti-depressant} 常指抑制 \textbf{serotonin/norepinephrine transporters} 的药物，这些转运体也是 Na$^+$ 依赖的。

\item \textbf{ATP-Driven Pumps有哪三种，分别都是在什么地方起作用}：三大类是 \textbf{P-type、ABC、F-type/V-type}。P-type 多见于质膜和 ER/SR 膜，转运离子如 Na$^+$、K$^+$、Ca$^{2+}$；ABC 广泛分布于细菌膜、真核质膜及细胞器膜，转运脂质、药物和代谢物；F-type 主要在线粒体内膜、叶绿体类囊体膜和细菌质膜上参与 ATP 合成，V-type 主要在 lysosome/endosome/vacuole 等酸化细胞器膜上泵 H$^+$。

\item \textbf{bacteria的膜内外转运的相互配合}：细菌常通过 \textbf{跨膜协同系统} 完成物质运输。典型如 ABC importer 与周质结合蛋白协作完成摄取；Gram-negative bacteria 的外排系统可由内膜转运体、周质连接因子和外膜通道共同组成，直接把底物跨两层膜排出细胞。

\item \textbf{F1F0 ATP synthase 的机理是什么？}：H$^+$ 顺电化学梯度经 \textbf{F$_0$} 回流，驱动转子旋转；旋转传递到 \textbf{F$_1$} 头部，使 $\beta$ 亚基在不同构象间循环，从而完成 ADP + P$_i$ 到 ATP 的合成。

\item \textbf{Bacterial Cells Against Extreme Osmotic Pressures依靠的是哪一种channel？}：依靠 \textbf{mechanosensitive channels}。

\item \textbf{钾离子通道蛋白的结构，以及它具有的三种状态？}：典型 K$^+$ channel 是 \textbf{四聚体}，每个亚基贡献跨膜螺旋和 \textbf{selectivity filter}，滤过区以羰基氧精确配位 K$^+$。若按门控状态概括，可分为 \textbf{closed、open、inactivated}。

\item \textbf{髓鞘的英文是什么？}：\textbf{myelin sheath}。

\item \textbf{乙酰胆碱受体的结构和机制}：神经肌接头的烟碱型乙酰胆碱受体是 \textbf{五聚体 ligand-gated cation channel}，由 5 条跨膜多肽组成。两分子 acetylcholine 结合后引发构象变化，通道开放，Na$^+$ 内流、K$^+$ 外流，导致终板去极化。

\item \textbf{Neuromuscular Transmission的激活过程（一直到t-tubule）}：动作电位到达运动神经末梢后，\textbf{电压门控 Ca$^{2+}$ 通道} 开放，促进 acetylcholine 释放；ACh 结合肌膜上的 \textbf{nicotinic ACh receptor}，引起终板去极化；达到阈值后，肌细胞膜上的 \textbf{voltage-gated Na$^+$ channels} 触发动作电位；动作电位沿肌膜传播并进入 \textbf{T-tubule}；T-tubule 上的电压感受器进一步耦联肌浆网的 Ca$^{2+}$ 释放通道，引发肌浆 Ca$^{2+}$ 升高和收缩。

\item \textbf{long-term potentiation的建立以及NMDA和AMPA的调控是怎样的？}：高频刺激先使突触后膜去极化，解除 \textbf{NMDA receptor} 上的 Mg$^{2+}$ 阻塞；谷氨酸存在时 NMDA 开放，Ca$^{2+}$ 内流，激活 CaMKII 等激酶；随后 \textbf{AMPA receptor} 被磷酸化并更多插入突触后膜，使同样刺激产生更大电流，形成 LTP。晚期 LTP 还伴随基因表达改变和树突棘结构增强。
\end{enumerate}

\addcontentsline{toc}{section}{细胞器与蛋白分选}
\section*{细胞器与蛋白分选}
\begin{enumerate}
\item \textbf{symbiotics的机制是什么，以及细胞核和内膜系统的形成？}：这里应指 \textbf{endosymbiotic theory}。线粒体和叶绿体被认为来源于古代原核生物被宿主吞入后长期共生；而 \textbf{细胞核与内膜系统} 更常解释为原始质膜向内折叠并逐渐特化形成的连续膜系统。

\item \textbf{蛋白质的三种不同的transport mechanisms，以及画出胞内转运机制以及方向图？}：三种机制是 \textbf{gated transport}、\textbf{transmembrane protein translocation}、\textbf{vesicular transport}。方向图如下：

\begin{verbatim}
Cytosol <-> Nucleus                : gated transport (NPC)
Cytosol -> ER                      : translocation
Cytosol -> Mitochondria/Chloroplast/Peroxisome : translocation
ER -> Golgi -> Plasma membrane     : vesicular transport
Golgi -> Endosome -> Lysosome      : vesicular transport
Plasma membrane -> Endosome        : endocytosis
Endosome -> Plasma membrane/Golgi  : recycling or retrograde transport
\end{verbatim}

\item \textbf{细胞核的内膜外膜}：核被 \textbf{inner nuclear membrane} 和 \textbf{outer nuclear membrane} 包围；外核膜与 ER 连续，内核膜富含与 chromatin 和 nuclear lamina 结合的特异蛋白。

\item \textbf{细胞核的import和export的具体机制}：核进出口经 \textbf{nuclear pore complex} 完成。带 \textbf{NLS} 的蛋白先被 importin 识别，再与核孔 FG-nucleoporins 弱相互作用穿过核孔；核内高 \textbf{Ran-GTP} 促使 cargo 释放。核输出相反，exportin 需同时结合 \textbf{NES cargo + Ran-GTP} 才能出核；到胞质后 RanGAP 促使 GTP 水解，复合物解离，实现方向性。

\item \textbf{线粒体转运涉及哪些蛋白，这些蛋白的功能分别是什么？}：\textbf{TOM} 是外膜总入口；\textbf{TIM23} 主要把前体蛋白送入基质或插入部分内膜；\textbf{TIM22} 主要插入多跨膜载体蛋白到内膜；\textbf{SAM} 负责外膜 $\beta$-barrel 蛋白装配；\textbf{OXA} 负责把某些基质侧蛋白再插入内膜；\textbf{mtHsp70} 等伴侣蛋白利用 ATP 帮助牵引和折叠；内膜电位也提供驱动力。

\item \textbf{蛋白质转运进类囊体需要几部分信号肽？}：通常需要 \textbf{两段信号}：先是导向叶绿体的 \textbf{transit peptide}，再是导向类囊体的第二段信号肽。

\item \textbf{过氧化物酶体的三种主要功能是什么？}：\textbf{超长链脂肪酸 $\beta$-氧化}、\textbf{plasmalogen 等特定脂质合成}、\textbf{氧化反应与 ROS/H$_2$O$_2$ 代谢解毒}。

\item \textbf{Adrenoleukodystrophy是什么？}：\textbf{肾上腺脑白质营养不良（ALD）} 是过氧化物酶体相关遗传病，因 VLCFA 代谢障碍导致超长链脂肪酸积累，损伤神经系统中的 \textbf{myelin sheath}。

\item \textbf{进入过氧化物酶体的信号肽在哪里？}：经典 PTS1 位于 \textbf{C 末端}，常见序列是 \textbf{Ser-Lys-Leu (SKL)} 或其变体。

\item \textbf{在肝脏的光滑内质网的功能是什么？}：肝细胞 smooth ER 主要负责 \textbf{脂质与脂蛋白合成}，并富含 \textbf{detoxification enzymes}，如 cytochrome P450。

\item \textbf{Rough ER和Smooth ER在结构上的区别是什么？}：\textbf{rough ER} 表面附着核糖体，多呈片层状；\textbf{smooth ER} 无核糖体，多呈管状网络。

\item \textbf{signal-recognition particle(SRP)的形成机制是什么？}：当新生肽 N 端 signal peptide 从核糖体伸出后，\textbf{SRP} 识别并结合该疏水信号肽及核糖体，短暂停止翻译；随后与 ER 膜上的 SRP receptor 对接，把核糖体送到 Sec61 translocon，翻译恢复并进行共翻译转运。

\item \textbf{signal peptide对于translocator的作用}：signal peptide 以 \textbf{短疏水 $\alpha$-helix} 形式插入 translocator，帮助 \textbf{打开 Sec61 通道}，并可充当 start-transfer signal。

\item \textbf{Post-translational translocation以及BiP在其中的作用？}：翻译完成后的蛋白可在伴侣帮助下送入 ER。Sec62/Sec63 等辅助复合体配合 Sec61 工作；ER 腔内 \textbf{BiP} 结合新出现的多肽段，利用 ATP 循环防止回滑，相当于把多肽“拉”入 ER 腔。

\item \textbf{形成single pass transmembrane protein的两次识别分别是什么？}：关键疏水信号序列先被 \textbf{SRP} 识别，再被 \textbf{protein translocator（Sec61 pore）} 识别。

\item \textbf{如何实现多跨膜区域？}：通过多组 \textbf{start-transfer} 和 \textbf{stop-transfer} 信号的组合，使多段疏水序列依次插入膜内，形成 multipass transmembrane protein。

\item \textbf{蛋白质的插入方向是如何选择的？}：主要遵循 \textbf{positive-inside rule}；即跨膜信号序列两侧正电荷分布不对称时，\textbf{正电荷更多的一侧更倾向留在胞质侧}。

\item \textbf{Tail-anchored Protein的锚定机制是什么？(有三个相关蛋白和一个complex)}：尾锚定蛋白在翻译后才暴露其 C 端跨膜段。常见哺乳动物框架中，\textbf{SGTA} 先捕获疏水尾段，\textbf{Bag6} 参与分选与质量控制，\textbf{TRC40}（ATPase）负责递送到底物受体；ER 膜上的 \textbf{WRB/CAML receptor complex} 接收并插入该蛋白。

\item \textbf{glycosylated的蛋白的修饰位点在哪里？}：若指 \textbf{N-linked glycosylation}，修饰加在 ER 腔内新生肽的 \textbf{Asn} 上，经典序列为 \textbf{Asn-X-Ser/Thr}（X 不为 Pro）。

\item \textbf{折叠失败的蛋白，不再尝试修复，而是被“拉回细胞质并降解”的过程？}：这是 \textbf{ER-associated degradation, ERAD}。

\item \textbf{IRE1机制}：ER stress 时，BiP 从 IRE1 上解离，IRE1 二聚/寡聚并自磷酸化，激活其 \textbf{RNase} 活性；随后剪接 \textbf{XBP1 mRNA}（酵母中是 HAC1），生成活性转录因子，上调分子伴侣、ERAD 和膜生成相关基因，从而增强 ER 折叠能力并缓解应激。
\end{enumerate}

\addcontentsline{toc}{section}{囊泡运输、Golgi、内吞与自噬}
\section*{囊泡运输、Golgi、内吞与自噬}
\begin{enumerate}
\item \textbf{vesicle function是哪三点？}：囊泡运输对 \textbf{细胞生长与分裂}、\textbf{内分泌和外分泌生理}、\textbf{神经生理活动} 都至关重要。

\item \textbf{clathrin蛋白包被的囊泡结构是什么？}：由 \textbf{clathrin triskelion} 组装成多面体外壳；其下层由 adaptor proteins 连接 cargo receptor 与 clathrin。

\item \textbf{COPII adaptor coat proteins是什么，他们有什么特点?}：核心 adaptor 是 \textbf{Sec23/Sec24}，构成 COPII \textbf{inner coat}；其中 Sec24 负责 cargo 识别，Sec23 与 Sar1 结合并参与 coat 装配。

\item \textbf{LIPID SIGNALS的机制是什么?}：主要指膜上的 \textbf{phosphoinositides} 作为空间标签，选择性招募带有特定 lipid-binding domain 的 coat/adaptor/effector 蛋白，从而规定囊泡生成位置和运输身份。

\item \textbf{recognizing the correct target membrane的两步骤？}：第一步是 \textbf{Rab + Rab effectors/tethering proteins} 介导初始 dock 到正确膜区；第二步是 \textbf{SNARE pairing} 促使两层膜真正融合。

\item \textbf{Rab的活跃和非活跃状态是由什么决定的？以及在活跃和非活跃状态下都有什么特点？}：由其 \textbf{结合 GTP 还是 GDP} 决定。Rab-GTP 为活跃态，牢固定位于膜上并招募 effectors；Rab-GDP 为非活跃态，常被 \textbf{GDI} 结合并保持在胞质中。GEF 激活它，GAP 使其失活。

\item \textbf{SNARE的两种分别有什么特点，他们的解离需要什么？}：\textbf{v-SNARE} 多位于囊泡膜，\textbf{t-SNARE} 多位于靶膜；两者形成极稳定的四螺旋复合体，提供膜融合能量。解离需要 \textbf{NSF ATPase} 和辅助因子 $\alpha$-SNAP，依赖 ATP 水解。

\item \textbf{ER到Golgi的正向运输和反向运输分别依靠的是什么包被蛋白？}：正向运输主要用 \textbf{COPII}；反向运输主要用 \textbf{COPI}。

\item \textbf{为什么需要retrieval pathway？}：因为部分 \textbf{ER resident proteins} 和运输机器会误入 Golgi，需要回收；此外也要把受体、SNARE 和其他循环利用组件送回起始站点，维持分区身份。

\item \textbf{Golgi的层级特点，大概每一层的功能？}：Golgi 由 \textbf{cis Golgi network, cis cisterna, medial cisterna, trans cisterna, trans Golgi network} 构成。CGN 负责接收来自 ER 的货物；cis 到 trans cisternae 依次完成糖链修饰与加工；TGN 负责最终分选，决定送往质膜、分泌颗粒或 lysosome/endosome。

\item \textbf{Golgi主要发生哪一种糖基化？}：Golgi 是 \textbf{O-linked glycosylation} 最重要的场所，同时也继续加工 N-linked oligosaccharides。

\item \textbf{蛋白质糖基化的意义？}：促进 \textbf{折叠与质量控制}、提高 \textbf{稳定性和抗蛋白酶性}、参与 \textbf{细胞识别与分选}、并在 \textbf{信号调控} 中发挥作用。

\item \textbf{怎么理解lysosome的hetergenous？}：lysosome 不是单一静态囊泡，而是与 late endosome、endo-lysosome 等不断融合分离、内容物和成熟阶段不同的一组异质性降解性细胞器。

\item \textbf{在植物和fungus里面什么承担了lysosome的功能？}：主要由 \textbf{vacuole} 承担类似 lysosome 的降解功能。

\item \textbf{在高尔基体里M6P pathway的标记和分选}：溶酶体水解酶在 \textbf{cis Golgi} 的 N-linked oligosaccharide 上被加上 \textbf{mannose-6-phosphate} 标记；到 \textbf{TGN} 后被 M6P receptor 识别并装入去往 endosome/lysosome 的囊泡。

\item \textbf{GlcNAc phosphotransferase标记过程是怎样的？}：该酶先识别 lysosomal hydrolase 的蛋白表面特征，把 \textbf{GlcNAc-phosphate} 转移到其 N-linked oligosaccharide 的 mannose 上；随后第二个酶切去外层 \textbf{GlcNAc}，暴露出 \textbf{M6P}。

\item \textbf{Endocytosis、Phagocytosis、Macropinocytosis都是什么？}：\textbf{endocytosis} 是从胞外摄入大分子和膜成分的总称；\textbf{phagocytosis} 是吞噬大颗粒或微生物，形成 phagosome；\textbf{macropinocytosis} 是通过膜褶皱非特异性大量摄入胞外液和膜。

\item \textbf{Autophagy的两种形式以及出现的原因和情况，吞噬过程中形成的结构是什么（一共五步）？}：两种主要形式是 \textbf{nonselective autophagy} 和 \textbf{selective autophagy}。前者常见于 \textbf{饥饿}，用于回收营养；后者用于清除 \textbf{受损线粒体、过氧化物酶体、核糖体、ER 或入侵病原体}。五步为：1. \textbf{Induction}；2. \textbf{isolation membrane/phagophore} 形成；3. 闭合成 \textbf{double-membrane autophagosome}；4. 与 lysosome 融合；5. 消化内膜与内容物。
\end{enumerate}

\addcontentsline{toc}{section}{补充：内吞、胆固醇与分泌}
\section*{补充：内吞、胆固醇与分泌}
\begin{enumerate}
\item \textbf{pinocytosis和phagocytosis的区别是什么？}：\textbf{pinocytosis} 摄取的是胞外液和小分子，颗粒小、较普遍；\textbf{phagocytosis} 摄取的是大颗粒、细菌或细胞碎片，通常由巨噬细胞、中性粒细胞等专职吞噬细胞完成，强依赖 actin 重塑。

\item \textbf{caveolae（小窝）介导的内吞路径是什么？}：caveolae 是富含 \textbf{cholesterol 和 sphingolipid} 的烧瓶状膜内陷，主要由 \textbf{caveolin/cavin} 稳定；内吞后可进入 caveolar vesicle，再与早期内体、回收内体或跨胞转运路径相连，也常作为信号平台。

\item \textbf{Macropinocytosis与其他方式的不同之处有哪些？}：它是 \textbf{非特异性、大体积、actin-driven} 的液体摄取方式；常由膜 ruffle 回折形成大囊泡，不依赖经典 clathrin coat，摄入体积明显大于一般受体介导内吞。

\item \textbf{atherosclerotic plaques的产生原因？}：常由 \textbf{LDL 在血管壁滞留并被氧化} 开始，随后被巨噬细胞经 scavenger receptors 吞入形成 \textbf{foam cells}；持续炎症、平滑肌迁移增殖和脂质沉积最终形成动脉粥样硬化斑块。

\item \textbf{胆固醇一般会存在在哪种particle里面？}：血液中的胆固醇主要存在于 \textbf{lipoprotein particles} 中，临床上最常强调的是 \textbf{LDL} 携带并向组织递送胆固醇，\textbf{HDL} 参与逆向转运。

\item \textbf{ESCRT Protein Complexes的主要功能是什么， 实现这个功能的主要机制是什么？}：ESCRT 主要负责 \textbf{reverse-topology membrane scission}，典型场景包括 multivesicular body 形成、病毒出芽、胞质分裂末端切断和膜修复。其核心机制是 \textbf{ESCRT-III} 在膜颈处聚合并收缩，\textbf{VPS4 ATPase} 再促使其重塑与解聚，完成切割。

\item \textbf{Recycling Endosomes Regulate Plasma Membrane Composition的一个最典型的例子是什么？}：经典例子是 \textbf{GLUT4} 在脂肪细胞和肌细胞中的循环。胰岛素刺激后，含 GLUT4 的囊泡从回收内体/储存库转运到质膜，提高葡萄糖摄取。

\item \textbf{Trans Golgi Network (TGN)出来有哪三条路径？}：三条典型路径是去 \textbf{质膜的 constitutive secretion}、去 \textbf{regulated secretory vesicles}、以及去 \textbf{endosome/lysosome} 的 M6P 路径。

\item \textbf{cargo concentration的机制是什么？}：依靠 \textbf{cargo receptors、adaptor proteins、coat assembly} 对货物进行选择性富集；某些可溶性分泌 cargo 还可在 TGN 腔内通过 \textbf{聚集/凝聚} 被浓缩。

\item \textbf{Inducible的分泌囊泡是如何实现效果的？（一共五步）}：可概括为 1. \textbf{在 TGN 分选并浓缩 cargo}；2. 形成 \textbf{immature secretory vesicle/granule}；3. 囊泡 \textbf{成熟}，包括酸化、蛋白水解和膜成分回收；4. 囊泡在质膜 \textbf{docking/priming}；5. 接到刺激后常由 \textbf{Ca$^{2+}$} 触发 SNARE 介导的快速胞吐。
\end{enumerate}

\addcontentsline{toc}{section}{补充：染色体与基因表达基础问答}
\section*{补充：染色体与基因表达基础问答}
\begin{enumerate}
\item \textbf{核纤层（nuclear lamina）的结构蛋白——lamins 的分类与类型是什么？}：nuclear lamina 的主要结构蛋白是 \textbf{lamins}，它们属于 \textbf{type V intermediate filaments}。按类型可分为 \textbf{A-type lamins}（如 lamin A、lamin C）和 \textbf{B-type lamins}（如 lamin B1、lamin B2）。

\item \textbf{核型分析（karyotyping）能检测哪些类型的染色体异常？}：可检测 \textbf{数目异常} 和 \textbf{较大的结构异常}，包括 \textbf{aneuploidy、trisomy、chromosomal translocation}，以及其他能在核型水平分辨的大片段重排。

\item \textbf{为什么DNA序列本身会影响核小体（nucleosome）的形成位置？}：因为不同 DNA 序列的 \textbf{弯曲能力不同}。核小体形成要求 DNA 紧密缠绕组蛋白核心，较易弯曲的序列更有利于核小体装配，因此 DNA sequence 会影响 nucleosome positioning。

\item \textbf{H1（连接组蛋白）的作用是什么？}：\textbf{linker histone H1} 结合核小体之间的 linker DNA，并改变 DNA 进出核小体的路径，从而 \textbf{促进染色质进一步压缩}。

\item \textbf{ATP依赖的染色质重塑复合体有什么作用？}：它们利用 \textbf{ATP hydrolysis} 不断 \textbf{滑动、拆卸或替换核小体}，从而调控 DNA 的可及性，使转录、复制和修复等过程所需蛋白能够接近 DNA。

\item \textbf{酵母着丝粒为什么说是由特殊核小体定义的？}：因为酵母着丝粒染色质中含有以 \textbf{CENP-A} 取代常规 H3 的特殊核小体，这种特化染色质定义 centromere 身份，并负责在分裂时建立与纺锤体的连接。

\item \textbf{异染色质（heterochromatin）为什么能self-propagating？}：因为其形成后可通过 \textbf{reader-writer 正反馈机制} 向邻近染色质扩散；reader 识别已有修饰，writer 在附近再写入同类修饰，因此异染色质状态能够传播并维持。

\item \textbf{研究染色体三维结构（3D genome）的方法有哪些？}：两大类常见方法是 \textbf{群体平均方法}，如 \textbf{3C/Hi-C}，用于测量群体细胞中的接触频率；以及 \textbf{单细胞成像方法}，如 \textbf{chromosome tracing}，用于直接观察单细胞中的空间构象。

\item \textbf{多线染色体（polytene chromosomes）的chromocenter是什么？}：多线染色体在 \textbf{着丝粒附近的异染色质区域} 会彼此连接并聚集，形成共同的 \textbf{chromocenter（染色中心）}。

\item \textbf{gene duplication 的意义是什么？}：\textbf{基因重复} 是产生新基因和新功能的重要来源之一。重复后，一个拷贝可保留原功能，另一个拷贝则更容易积累变异，从而演化出新功能或新调控方式。

\item \textbf{真核细胞的三种RNA聚合酶（Pol I / II / III）分别转录什么？}：\textbf{Pol I} 主要转录 rRNA；\textbf{Pol II} 主要转录 mRNA 以及部分 noncoding RNA；\textbf{Pol III} 主要转录 tRNA、5S rRNA 和其他小 RNA。

\item \textbf{topoisomerase 为什么重要？}：DNA 在复制或转录时会产生 \textbf{superhelical tension}；\textbf{topoisomerase} 负责快速切开并重新连接 DNA，从而消除扭转应力，保证复制和转录顺利进行。

\item \textbf{真核mRNA最先发生的加工步骤是什么？}：最先发生的是 \textbf{5' capping}。当新生 RNA 刚转录出约 25 个核苷酸时，其 5' 端便被加上帽结构。

\item \textbf{真核mRNA的3'端是如何形成的？}：真核 mRNA 的 3' 端通常不是直接转录终止得到，而是通过 \textbf{先切割，再加 poly(A) tail} 的方式加工形成。

\item \textbf{DNA looping 在细菌转录激活中的意义是什么？}：DNA 可以弯曲成环，使远端调控蛋白与启动子附近的 RNA polymerase 接触，从而 \textbf{增强或启动转录}。

\item \textbf{什么是Transcriptional synergy？}：指多个转录激活因子共同作用时产生 \textbf{大于简单相加} 的转录增强效应，即“\textbf{1+1 >> 2}”。

\item \textbf{Insulator（绝缘子）DNA序列的作用是什么？}：insulator 像染色质上的 \textbf{边界/隔离墙}，可以阻断 enhancer 等远程调控元件对不应被影响基因的作用，从而限制调控范围。

\item \textbf{什么是cell memory（细胞记忆）？}：短暂信号可通过 \textbf{positive feedback} 或稳定转录回路，被转化为长期持续的基因表达状态；即使最初信号消失，细胞仍能维持既定身份。

\item \textbf{Genomic imprinting是什么？}：genomic imprinting 是一种 \textbf{亲本来源依赖的单等位基因表达} 现象，即同一基因在父源和母源等位基因中只有一方表达，主要依赖表观遗传修饰。

\item \textbf{X 染色体失活（XCI）是什么？}：在哺乳动物雌性细胞中，一条 X 染色体会被 \textbf{大范围沉默}，以实现 dosage compensation。这个过程与 \textbf{Xist lncRNA} 和染色质沉默密切相关。

\item \textbf{同一个基因如何通过APA决定蛋白是膜型还是分泌型？}：通过选择不同的 \textbf{alternative cleavage and polyadenylation（APA）} 位点，可产生不同 3' 端的 mRNA；若较上游位点使编码跨膜区的下游序列不再保留，则更容易生成 \textbf{分泌型} 蛋白，反之可保留 \textbf{膜型} 蛋白。
\end{enumerate}

\addcontentsline{toc}{section}{第二部分：灵活问答}
\section*{第二部分：灵活问答}
本部分延续前文格式，但题目改为更偏向“意义、机制、比较、应用和因果推断”的问法，尽量覆盖笔记中已经出现的知识点。

\addcontentsline{toc}{subsection}{细胞与模型生物}
\subsection*{细胞与模型生物}
\begin{enumerate}
\item \textbf{为什么说细胞是生命的基本单位，这句话在实验生物学里真正意味着什么？}：它意味着所有生物系统都建立在细胞这一共同组织层次上，很多看似复杂的生命现象都可以拆解为细胞层面的结构、代谢、信息储存和分裂行为来研究，因此用 model organism 得到的规律常常可以推广到更高等生物。

\item \textbf{为什么rRNA序列特别适合用来构建“生命之树”？}：因为 rRNA 在所有细胞中都存在、功能核心而高度保守，同时又保留了足够的变异位点，既能比较远缘物种，也能区分近缘物种，所以非常适合做系统发育分析。

\item \textbf{为什么说古菌在分子层面上往往更接近真核生物，而不是更接近细菌？}：因为古菌在 DNA 复制、转录和染色质组织等很多信息处理系统上与真核更相似，说明“原核外形相似”不等于“分子系统相近”。

\item \textbf{为什么不同细胞能有完全不同的形态和功能，但仍然拥有相同的DNA？}：关键在于 \textbf{spatiotemporal gene expression} 不同。细胞通过选择性转录、RNA 加工、翻译和蛋白稳定性控制，建立不同的蛋白组合，因此在相同基因组背景下形成不同命运。

\item \textbf{为什么模式生物在细胞生物学里如此重要？}：因为它们通常生命周期短、易培养、遗传工具丰富、成本低、伦理负担小，而且很多核心细胞机制高度保守，因此能快速建立可重复的因果实验。

\item \textbf{为什么酵母既适合研究细胞周期，也适合研究自噬？}：酵母是简单真核，既保留 nucleus、organelle、cytoskeleton 等核心结构，又易做突变和筛选；细胞周期表型容易观察，自噬则可通过营养饥饿和液泡降解阻断直接积累中间结构。
\end{enumerate}

\addcontentsline{toc}{subsection}{显微与实验技术}
\subsection*{显微与实验技术}
\begin{enumerate}
\item \textbf{为什么分辨率比放大倍数更能决定显微镜“有用不好用”？}：因为放大只会把图像变大，不能凭空增加细节；只有分辨率决定两个相邻结构能否被区分，所以它才是显微技术最核心的性能指标。

\item \textbf{相差显微镜为什么特别适合观察活细胞？}：因为它把透明样品造成的相位差转换成亮暗差，不需要染色即可观察活细胞内部结构，因此能减少损伤并保留动态过程。

\item \textbf{荧光显微镜相比普通光学显微镜最大的概念提升是什么？}：不是单纯“更亮”，而是引入了 \textbf{molecular specificity}，即可以只看某一类分子或某一结构，从“看到细胞”变成“看到特定组分在细胞中的位置和动态”。

\item \textbf{FRAP实验最适合回答什么类型的问题？}：最适合回答蛋白或脂质在膜中的 \textbf{mobility}、扩散速度和是否受到局部约束的问题，而不是直接告诉你分子总量或结合位点。

\item \textbf{为什么细胞融合实验能证明膜蛋白会在膜平面内扩散？}：不同荧光标记的膜蛋白在异核细胞融合后会逐渐混合，说明它们不是固定在原地，而是能在脂双层中侧向扩散。

\item \textbf{为什么SDS-PAGE能按分子量分离蛋白？}：因为 SDS 让蛋白基本都带上与长度近似成比例的负电荷，并破坏原有折叠差异，于是迁移率主要由分子大小决定。

\item \textbf{二维电泳相比普通SDS-PAGE多解决了什么问题？}：它先按等电点分，再按分子量分，因此能把“大小接近但电荷不同”的蛋白拆开，提高复杂样本解析度。

\item \textbf{为什么X射线晶体学和冷冻电镜都属于“结构生物学”，但面对的问题并不完全相同？}：因为它们都提供高分辨率结构，但晶体学依赖结晶，更适合稳定样品；冷冻电镜对大复合体和难结晶样品更友好，更适合看天然构象异质性。
\end{enumerate}

\addcontentsline{toc}{subsection}{蛋白质、DNA与基因表达方法}
\subsection*{蛋白质、DNA与基因表达方法}
\begin{enumerate}
\item \textbf{为什么“保守序列”常常能提示功能重要性？}：因为功能关键位点受到 purifying selection，突变更不容易被长期保留下来；跨物种保守往往说明该序列在结构、催化或调控上有重要作用。

\item \textbf{为什么基因敲除、过表达和定位实验往往要结合起来看？}：因为“缺失后会怎样”告诉你必要性，“增加后会怎样”提示剂量和调控，“定位在哪里”提供作用场景，三者结合才更容易建立机制。

\item \textbf{为什么RNA水平和蛋白水平不一定一致？}：因为从 DNA 到蛋白中间还有 RNA 加工、输出、翻译效率和蛋白稳定性等多层调控，所以 mRNA abundance 只是蛋白表达的上游指标，不是最终结果。

\item \textbf{为什么替代剪接能显著提高真核细胞的功能复杂性？}：因为一个基因可以生成多个不同 mRNA 与蛋白异构体，从而在不显著增加基因数目的情况下扩展功能库。

\item \textbf{为什么microRNA更像“微调器”而不是简单的开关？}：因为一个 miRNA 往往弱到中等程度地调控一组 mRNA，而不是完全开或关一个基因，因此更像网络级别的精细缓冲和协调器。
\end{enumerate}

\addcontentsline{toc}{subsection}{细胞核、染色体与基因组}
\subsection*{细胞核、染色体与基因组}
\begin{enumerate}
\item \textbf{为什么核仁通常被视为最典型的nuclear body？}：因为它是最明显、最大、功能最集中的核内无膜区室，主要承担 rRNA 合成和核糖体亚基装配。

\item \textbf{核纤层除了“支撑细胞核”之外还有什么更深层的意义？}：它不仅提供机械支撑，还参与染色质定位、核孔锚定、DNA 复制和有丝分裂相关过程，因此既是结构框架也是调控平台。

\item \textbf{为什么说DNA双螺旋结构天然就暗示了遗传复制机制？}：因为两条链互补，一条链就可以作为另一条链复制时的模板，所以结构本身就解释了遗传信息如何被忠实传递。

\item \textbf{为什么真核DNA必须包装成染色体，而不能“裸奔”？}：因为真核基因组极长，若不压缩就无法装入细胞核；同时包装还能保护 DNA，并通过染色质状态调节复制、转录和修复。

\item \textbf{为什么nucleosome既帮助压缩DNA，又会给转录带来阻碍？}：因为 DNA 绕在组蛋白八聚体上后更紧凑、更稳定，但同时也降低了转录因子和 RNA polymerase 对 DNA 的可及性。

\item \textbf{ATP-dependent chromatin remodeling complex 的真正作用逻辑是什么？}：它们不是“打开所有染色质”，而是通过滑动、移除或替换核小体，\textbf{局部地} 改变 DNA 可及性，让特定反应在特定位置发生。

\item \textbf{为什么histone modification 常常被说成是一种“语言”或“代码”？}：因为不同修饰不仅改变组蛋白本身电荷和物理性质，还能招募特定读取蛋白，组合起来形成可被识别和传播的调控状态。

\item \textbf{为什么heterochromatin 能自我传播？}：因为很多抑制性修饰体系同时包含 \textbf{reader 和 writer}，reader 识别已存在修饰后招募 writer 到邻近核小体，从而把同样修饰向外扩展。

\item \textbf{为什么centromere 的身份不能只用DNA序列来解释？}：因为尤其在高等真核中，centromere 还依赖 \textbf{CENP-A} 等特殊染色质状态维持，说明它带有明显的表观遗传成分。

\item \textbf{为什么说染色体高级结构是“fluid and hierarchical”的？}：因为它既有 loops、TADs、compartments 和 territories 等多层级组织，又会随细胞类型、转录状态和细胞周期动态变化。

\item \textbf{polytene chromosome 为什么对研究染色体结构特别有用？}：因为它含有成百上千份平行对齐的 DNA，带纹和 puff 都很清楚，能直接在光镜下观察局部解凝缩与基因表达。

\item \textbf{为什么基因表达活跃区域常会从染色体territory中伸出来？}：因为高转录活性常伴随局部染色质去凝缩，使该区域更易接触转录机器和调控因子。
\end{enumerate}

\addcontentsline{toc}{subsection}{转录与转录调控}
\subsection*{转录与转录调控}
\begin{enumerate}
\item \textbf{RNA相对DNA最重要的化学差异为什么会带来巨大的结构差异？}：因为 RNA 多了 2'-OH，而且通常是单链，这使它更容易形成复杂二级和三级结构，因此既能携带信息，也能承担结构和催化功能。

\item \textbf{为什么细菌和真核生物的转录起始复杂度差异这么大？}：因为真核 DNA 被包装进染色质，启动子可及性低、调控距离更远，所以必须借助更多一般转录因子、调控蛋白和染色质调节机制。

\item \textbf{为什么RNA polymerase II 的CTD如此重要？}：因为 CTD 像一个可编程平台，不同磷酸化状态可以顺序招募 capping、splicing、3' end processing 等机器，把转录和 RNA 加工耦联起来。

\item \textbf{为什么真核转录常常和RNA加工同步发生？}：这样能提高效率并减少错误，把新生 RNA 一边合成一边加帽、剪接和加 poly(A)，更快生成可输出和可翻译的成熟 mRNA。

\item \textbf{为什么alternative splicing 是细胞特异性的重要来源？}：因为不同细胞可以用不同 splicing regulators 选择不同外显子组合，从而制造细胞类型特异的蛋白功能。

\item \textbf{为什么eukaryotic gene regulation 常常强调 combinatorial control？}：因为单个 regulator 通常不足以决定细胞命运，真正决定转录输出的是多个 transcription factors、chromatin state 和 enhancer 组合后的整体逻辑。

\item \textbf{为什么转录激活常常是“协同”的，而不是简单相加？}：因为多个 activator 可共同招募 coactivator、染色质重塑复合体和 RNA polymerase machinery，一个因子提高另一个因子的结合和作用效率，因此效果常是 1+1 远大于 2。

\item \textbf{为什么细胞记忆能在没有DNA序列改变的情况下长期存在？}：因为转录正反馈回路、DNA methylation 和染色质修饰状态都能在细胞分裂中延续，使同一谱系细胞长期保持既定身份。

\item \textbf{X染色体失活为什么是研究表观遗传的经典模型？}：因为它展示了 lncRNA、染色质沉默和细胞记忆如何协同，让一整条染色体长期转录沉默并在后代细胞中稳定传递。
\end{enumerate}

\addcontentsline{toc}{subsection}{核、染色体与转录调控补充机制题}
\subsection*{核、染色体与转录调控补充机制题}
\begin{enumerate}
\item \textbf{为什么nuclear lamina 的改变不仅会影响细胞核形状，还会影响基因表达？}：因为 nuclear lamina 一方面提供机械支撑，另一方面还是染色质和核孔的锚定位点；因此 lamin 异常会同时破坏 \textbf{核结构、染色质组织和基因调控}。

\item \textbf{为什么karyotyping 适合发现大片段异常，却不适合发现点突变？}：因为核型分析观察的是整条染色体的数目和宏观带型，它擅长看 \textbf{整条染色体增减或大片段重排}，但分辨率远不足以识别 nucleotide-level 变化。

\item \textbf{为什么说核小体定位既受DNA序列影响，又不能完全由DNA序列决定？}：因为某些序列更容易弯曲，先天地更适合缠绕组蛋白；但细胞内还有 H1、转录因子和 ATP-dependent remodeling complexes 持续移动核小体，所以最终定位是 \textbf{序列偏好和调控活动共同决定} 的。

\item \textbf{H1 为什么会让染色质更“高级压缩”，而不只是多放一个组蛋白？}：因为 H1 结合 linker DNA 并改变 DNA 离开核小体的角度，使相邻核小体更容易靠近和堆积，因此直接促进更高阶染色质纤维形成。

\item \textbf{为什么染色质重塑复合体对转录调控特别关键？}：因为转录调控的瓶颈常不是 DNA 有没有序列，而是 DNA \textbf{能不能被接触到}；重塑复合体通过挪开或重排核小体，把“有信息的 DNA”变成“可被读取的 DNA”。

\item \textbf{为什么CENP-A 核小体可以被理解为着丝粒身份的“表观遗传标记”？}：因为它让着丝粒的定义不完全依赖普通 DNA 序列，而依赖一种可在细胞代际间延续的特殊核小体状态，从而稳定指定纺锤体连接位点。

\item \textbf{为什么heterochromatin 的reader-writer 机制会天然导致“边界扩散”的风险？}：因为一旦 reader 识别到已有抑制标记，它就会帮助 writer 把同样标记写到邻近区域；若没有边界限制，这种正反馈就可能持续向外蔓延，误伤本应活跃的基因。

\item \textbf{为什么研究3D genome 需要同时用Hi-C 和 chromosome tracing 这两类方法？}：因为 Hi-C/3C 给出的是 \textbf{群体平均接触频率图谱}，善于概括全局规律；chromosome tracing 则直接显示 \textbf{单细胞空间构象}，善于揭示细胞间异质性。两者互补，才能避免把平均图误当成每个细胞都一样。

\item \textbf{为什么polytene chromosome 的chromocenter 有助于理解异染色质的组织作用？}：因为它清楚展示了着丝粒附近异染色质不仅是“沉默区域”，还承担 \textbf{把多条染色体聚拢组织起来} 的结构作用。

\item \textbf{为什么gene duplication 被看作演化创新的安全机制？}：因为复制后一个拷贝仍可维持原功能，另一个拷贝获得了“试错空间”，更容易积累突变并发生 \textbf{subfunctionalization 或 neofunctionalization}。

\item \textbf{为什么真核细胞需要三种RNA polymerase，而不是一个polymerase做完所有事？}：因为不同 RNA 类型在长度、结构、启动子要求、加工方式和表达规模上差异很大，使用专门化 polymerase 能提高 \textbf{效率、准确性和调控灵活性}。

\item \textbf{为什么topoisomerase 在高转录活性区域尤其重要？}：因为 RNA polymerase 沿 DNA 前进会持续制造正负超螺旋，若不及时释放，DNA 会变得难以展开，进而反过来阻碍后续复制或转录。

\item \textbf{为什么5' capping 会在转录刚开始不久就发生，而不是等mRNA做完再统一处理？}：因为帽结构需要尽早保护 RNA 5' 端免受降解，并为后续剪接、核输出和翻译起始提供识别标记，所以它是最早发生的加工步骤。

\item \textbf{为什么真核mRNA 的3'端形成要靠“切割+poly(A)”而不是直接照着DNA模板转完？}：因为这样更利于调控转录终止、mRNA 稳定性、核输出和翻译效率，而且同一基因还能通过改变切割位点产生不同转录本。

\item \textbf{为什么DNA looping 能显著增强转录调控的灵活性？}：因为它让线性序列上相距很远的调控元件在三维空间中彼此接触，从而打破“只能邻近调控”的限制。

\item \textbf{为什么transcriptional synergy 经常意味着多个调控因子不是独立工作的？}：因为它提示这些因子之间存在协作招募、稳定彼此结合或共同改造局部染色质环境的过程，所以总效应远超简单相加。

\item \textbf{为什么insulator 对复杂基因组尤其重要？}：因为真核基因组中 enhancer 很多且作用距离远，如果没有 insulator 设边界，远程激活信号就容易跨域影响错误靶基因。

\item \textbf{为什么positive feedback 是产生cell memory 的经典回路结构？}：因为一旦某个 regulator 被短暂信号推高，它就能继续促进自身或同伴表达，即使外界信号消失，系统也能维持在“开”的状态。

\item \textbf{为什么genomic imprinting 会被认为是“谁给的等位基因很重要”？}：因为父源和母源等位基因并不等价，细胞会通过甲基化等标记记住其来源，并选择只表达其中一方。

\item \textbf{为什么XCI 是 dosage compensation，而不是简单关掉一条染色体？}：因为它的本质目的是让 XX 个体与 XY 个体在 X-linked gene dosage 上保持平衡，所以它是一种全基因组剂量校正机制。

\item \textbf{为什么APA 可以改变蛋白最终去向，而不仅仅是改变mRNA长度？}：因为不同 poly(A) 位点可能决定下游编码区是否保留；一旦跨膜区编码序列被提前截断，同一个基因就能从 \textbf{membrane-bound form} 切换为 \textbf{secreted form}。
\end{enumerate}

\addcontentsline{toc}{subsection}{细胞骨架机制题}
\subsection*{细胞骨架机制题}
\begin{enumerate}
\item \textbf{为什么actin filament 特别适合驱动细胞边缘快速形变，而微管不适合承担这个角色？}：因为 actin 更细、更灵活、聚合去聚合更快，而且主要位于皮层和质膜下方，适合推动 lamellipodia、filopodia 和收缩环；微管更适合长距离运输和建立空间坐标。

\item \textbf{为什么actin nucleation 会成为聚合反应的限速步骤？}：因为一开始少数单体形成稳定 nucleus 在热力学上不容易，只有跨过这个初始门槛后，后续 elongation 才会快速进行。

\item \textbf{为什么ATP hydrolysis 会让actin filament 出现 treadmilling？}：因为 plus end 和 minus end 的临界浓度不同，当游离单体浓度处于两者之间时，正端净装配、负端净去装配，丝长不变但亚基持续“流动”。

\item \textbf{formin 和 Arp2/3 的分工为什么能直接决定细胞边缘长什么样？}：因为 formin 生成长而直的丝，更适合做 bundle；Arp2/3 在已有丝侧面分支，更适合形成致密 branched network，所以前者更像“搭梁”，后者更像“织网”。

\item \textbf{为什么肌球蛋白II必须形成bipolar thick filament 才特别适合做收缩机器？}：因为双极结构能让两侧头部朝相反方向沿 actin 走，从而把反向极性的 actin filament 往中间拉，产生整体收缩而不是单向运输。

\item \textbf{为什么微管的动态不稳定性对有丝分裂尤其关键？}：因为纺锤体必须不断试探、捕获并重新组织染色体连接点，若微管过于稳定或过于不稳定，都无法高效完成染色体排列和分离。

\item \textbf{为什么中间丝被称为“机械韧性系统”，而不是“运输系统”？}：因为它们不像 actin 和微管那样承担明显方向性的 motor-based transport，而是更擅长承受拉伸、分散应力、维持组织完整性。

\item \textbf{为什么上皮组织中特别依赖中间丝和连接结构的整合？}：因为上皮要长期承受剪切和张力，keratin 与 desmosome/hemidesmosome 相连后能把单细胞的受力整合为组织级抗拉网络。
\end{enumerate}

\addcontentsline{toc}{subsection}{生物膜与电生理机制题}
\subsection*{生物膜与电生理机制题}
\begin{enumerate}
\item \textbf{为什么说膜的“流动性”不是一个装饰性性质，而是决定功能的核心参数？}：因为膜要完成运输、信号、分选、内吞、胞吐和细胞迁移，都要求脂质和蛋白具有一定流动性；但它又不能过于松散，否则屏障性和分区性会下降。

\item \textbf{为什么脂双层本身就足以形成屏障，但远不足以形成“活的膜”？}：因为纯脂双层只提供被动隔离，而真正的生物膜还需要蛋白来执行选择性转运、信号识别、黏附和酶促反应。

\item \textbf{为什么通道比转运体快得多，但在选择性控制上反而更依赖“门控”？}：因为通道一旦打开就允许大量离子快速流过，所以必须通过电压、配体或机械力门控来严格控制何时开放。

\item \textbf{为什么Na$^+$/K$^+$ pump 对神经兴奋性如此重要，即使它不直接产生动作电位？}：因为它长期建立 Na$^+$ 和 K$^+$ 梯度，这是静息电位和动作电位能够存在的前提；没有它，离子梯度会逐渐耗散，兴奋性最终消失。

\item \textbf{为什么K$^+$ channel 能高效区分K$^+$和更小的Na$^+$？}：因为 selectivity filter 提供的是适合 \textbf{脱水后的 K$^+$} 的精确配位几何，而 Na$^+$ 太小，无法被这些羰基氧稳定补偿脱水代价，所以反而不容易通过。

\item \textbf{为什么myelination 会显著提高神经传导速度？}：因为髓鞘减少离子泄漏并提高膜电阻，使去极化电流能传播更远，动作电位主要在 Ranvier 结间跳跃式再生。

\item \textbf{为什么NMDA receptor 能充当“巧合检测器”？}：因为它既需要谷氨酸结合，又需要膜去极化解除 Mg$^{2+}$ 阻塞，只有突触前释放和突触后激活同时发生时才大量通 Ca$^{2+}$。
\end{enumerate}

\addcontentsline{toc}{subsection}{蛋白分选与细胞器机制题}
\subsection*{蛋白分选与细胞器机制题}
\begin{enumerate}
\item \textbf{为什么所有蛋白都在胞质核糖体上起始翻译，却还能被送往完全不同的细胞器？}：因为蛋白自身携带 \textbf{sorting signal}，并被不同受体和转运机器识别，决定其后续命运。

\item \textbf{为什么核孔运输和线粒体/ER转运在概念上必须区分开？}：因为核孔允许已折叠蛋白通过，并通过 Ran 系统赋予方向性；而 ER、线粒体等跨膜转运通常要求多肽至少部分展开，是真正意义上的穿膜过程。

\item \textbf{为什么线粒体导入常常需要膜电位？}：因为很多线粒体前导肽带正电，内膜电位可电驱动其朝基质侧移动，同时与 ATP 依赖的伴侣系统协同完成导入。

\item \textbf{为什么过氧化物酶体在代谢病中经常与神经系统病变联系在一起？}：因为它不仅分解 VLCFA，还参与 plasmalogen 合成；这类脂质对 myelin 很重要，所以其缺陷常首先伤及神经系统。

\item \textbf{为什么粗面内质网常被视为“分泌途径的入口站”？}：因为分泌蛋白、溶酶体蛋白和多数膜蛋白都在这里进入内膜系统，并在此开始折叠、组装和 N-糖基化。

\item \textbf{为什么SRP 的存在能提高蛋白分选准确性？}：因为它在翻译早期就捕获 signal peptide，并把核糖体直接送到 ER 膜，避免疏水分泌蛋白在胞质中错误折叠或聚集。

\item \textbf{为什么BiP 既是折叠伴侣，也是转运机器的一部分？}：因为它不仅帮助蛋白正确折叠，还在 post-translational translocation 中结合新进入 ER 腔的多肽，防止回滑并推动单向导入。

\item \textbf{为什么ER质量控制必须同时具备“再尝试折叠”和“最终降解”两套机制？}：因为很多蛋白只是暂时未成熟，应该被挽救；但若持续错误折叠，则必须经 ERAD 移除，否则会引发蛋白毒性和 ER stress。

\item \textbf{为什么UPR 不能只做一件事，而要同时“增加折叠能力”和“减少新负担”？}：因为单纯增加伴侣不一定够，若合成压力继续过高，ER 仍会失衡；所以它必须同时扩容和限流。
\end{enumerate}

\addcontentsline{toc}{subsection}{囊泡运输与膜交通机制题}
\subsection*{囊泡运输与膜交通机制题}
\begin{enumerate}
\item \textbf{为什么vesicle transport 常被说成是“保持细胞区室身份”的系统，而不只是送货系统？}：因为囊泡运输不仅搬运 cargo，也持续回收和重新分配膜蛋白、SNARE、受体与脂质，维持各细胞器不同的组成和功能。

\item \textbf{为什么coat protein 不能单独决定囊泡命运？}：因为 coat 主要负责出芽和初始货物选择，而最终靶向还依赖 Rab、tether 和 SNARE，所以 coat 只是运输身份的一部分。

\item \textbf{为什么phosphoinositide 可以被理解为“膜地址码”？}：因为不同膜区富集不同磷酸化形式的 PI，它们能特异招募对应效应蛋白，让相同脂双层出现不同“功能地标”。

\item \textbf{为什么Rab和SNARE要分成两步，而不让SNARE直接完成所有特异性识别？}：因为 Rab/tether 负责远距离、可逆、初筛式的定位；SNARE 则负责近距离、强耦合、不可逆的真正融合，两步分工可提高准确性并降低误融合。

\item \textbf{为什么ER需要retrieval pathway，而不能简单地“只让该出去的出去”？}：因为分选永远不是百分之百准确，且许多运输元件本来就要循环使用，所以必须有回收机制来纠偏并维持 ER 身份。

\item \textbf{为什么Golgi 特别适合作为糖基化加工中心？}：因为它具有层级分明的 cisternae，每层有不同 glycosidase 和 glycosyltransferase，货物经过时可按顺序逐步修饰。

\item \textbf{为什么lysosome 被称为异质性细胞器而不是单一囊泡？}：因为其内容物、成熟阶段、融合历史和膜组成持续变化，与 late endosome 和 endo-lysosome 处于动态互变之中。

\item \textbf{为什么M6P 标记必须在Golgi 里而不是在溶酶体里完成？}：因为它的作用是 \textbf{提前分选} 新合成的 lysosomal hydrolases，必须在它们被送往 lysosome 之前就加上并被受体识别。

\item \textbf{为什么自噬既是“生存机制”，又可能与细胞死亡相关？}：因为适度自噬能回收营养、清除受损细胞器、缓解压力；但若长期过强或与其他损伤叠加，也可能反映或促进不可逆细胞衰败。

\item \textbf{为什么选择性自噬比非选择性自噬更能体现细胞“主动管理质量”的能力？}：因为它不是随机吞一团胞质，而是能识别特定受损目标，如 mitochondria、peroxisome 或 microbes，说明细胞具备精确清除系统。

\item \textbf{为什么受体介导内吞和macropinocytosis 在生理意义上非常不同？}：前者突出 \textbf{specificity and efficiency}，适合回收特定配体；后者突出 \textbf{bulk uptake}，适合大规模摄取液体、膜和外界物质。

\item \textbf{为什么ESCRT 在膜生物学里非常特别？}：因为它负责的是与 clathrin/COPI/COPII 出芽相反方向的膜切割，即从膜颈内部向外完成 \textbf{reverse-topology scission}。
\end{enumerate}

\end{document}
